\chapter{The discretised Furstenberg estimate}
\label{chap:furstenberg}

The radial projection argument, and through it the $C^2$ grain
argument, uses the incidence gain of Orponen and Shmerkin~\cite{OS}.
We use their September 2023 arXiv version. Its regular-case theorem
is Theorem 6.2 and its scale-combining proposition is Proposition 7.5;
the project and some earlier numbering call these Theorem 6.1 and
Proposition 7.3. Proposition 8.1 retains the same number.

\section{Point and tube configurations}

\begin{definition}[Discrete nonconcentration; OS, Definitions 2.1 and 2.8]
\label{def:os-set}
\lean{IsDeltaSSet}
In a metric space, a nonempty bounded set $P$ is a
$(\delta,s,C)$-set when
\[
    \cE_\delta(P\cap B(x,r))\leq Cr^s\cE_\delta(P)
    \quad(r\geq\delta).
\]
For planar lines, use a metric on the affine Grassmannian; in a
bounded nonvertical chart this is comparable to the metric on
slope-intercept parameters.
\end{definition}

\begin{definition}[Dyadic tubes; OS, Definitions 2.9--2.10]
\label{def:dyadic-tube}
The point-line duality sends $(a,b)$ to the line $y=ax+b$.
A dyadic $\delta$-tube is the union of all such lines whose parameter
pair lies in a dyadic square of side $\delta$, with slopes in a
fixed bounded chart. Its slope is the left endpoint of the
corresponding slope interval. Within a bounded spatial square it
is comparable to a usual tube of width $\delta$.
\end{definition}

\begin{definition}[Nice configurations; OS, Definition 5.1]
\label{def:nice-configuration}
\uses{def:os-set,def:dyadic-tube}
A $(\delta,s,C,M)$-nice configuration consists of a set $\cP$ of
dyadic $\delta$-squares and, for every $p\in\cP$, a
$(\delta,s,C)$-family $\T(p)$ of exactly $M\geq1$ dyadic tubes
meeting $p$. Write $\T=\bigcup_p\T(p)$. No distribution condition
on $\cP$ is included in this definition.
\end{definition}

\begin{theorem}[Discretised Furstenberg estimate; OS, Theorem 1.3]
\label{thm:os-discretized}
\lean{discretised_furstenberg_estimate}
\leanok
\uses{def:os-set}
For $0<s<1$ and $s<t<2$, there is $\eps(s,t)>0$ such that the
following holds at sufficiently small scales $\delta$. Let
$X\subset B(0,1)\subset\R^2$ be a
$(\delta,t,\delta^{-\eps})$-set. For every $x\in X$, let
$\T(x)$ be a $(\delta,s,\delta^{-\eps})$-set of lines whose
$\delta$-neighbourhoods contain $x$. Then
\[
                 \cE_\delta\left(\bigcup_{x\in X}\T(x)\right)
                             \geq\delta^{-2s-\eps}.
\]
The function $\eps(s,t)$ can be chosen bounded away from zero on
each compact subset of the open parameter range.
\end{theorem}
\begin{proof}
\uses{thm:os-regular,prop:os-combining,prop:os-multiscale,
  lem:os-uniformization}
\leanok
Pass to a uniform dyadic model. Decompose its scales into structured
and bad intervals. Structured intervals of dimension near $s$ incur
no substantial loss, while a positive proportion of the logarithmic
scale range has dimension strictly larger than $s$. On that portion
the regular-case incidence theorem gives a gain. Combine the
estimates across scales and choose the bad-scale and uniformization
losses smaller than that gain. Small perturbations of $(s,t)$ allow
one choice of gain on a neighbourhood; compactness then gives the
last assertion.
\end{proof}

\begin{remark}[The range and the geometric model]
\label{rem:os-range}
The paper includes $t=2$ and first states the theorem for dyadic
squares and tubes. The exported Lean theorem above uses ordinary
line neighbourhoods and the open range $s<t<2$, including compact
uniformity in that range. These are the parameters needed by the
subsequent arguments. The passage between models includes an actual
incidence-preserving enlargement: rounding a line's parameters alone
does not guarantee that the rounded line meets the same point cell.
\end{remark}

\begin{theorem}[Bounded charts; OS, Theorem 3.1]
\label{thm:os-bounded}
\uses{def:os-set,def:dyadic-tube}
Fix a bounded chart of radius $R$. The dyadic form of
Theorem~\ref{thm:os-discretized} holds for point squares in
$[-R,R)^2$ and tubes in the corresponding slope chart, with the
scale threshold also depending on $R$. Incidence may be imposed
between the closures of the tube and the square. The conclusion is
$\#\T\geq\delta^{-2s-\eps}$.
\end{theorem}
\begin{proof}
\uses{thm:os-discretized}
Partition the bounded chart into a fixed number of patches, rescale
one substantial patch, and allow a fixed enlargement of the tubes
and point cells. Absorb the resulting constants in the exponent loss.
\end{proof}

\begin{theorem}[Dual form; OS, Theorem 3.2]
\label{thm:os-dual}
\uses{def:os-set,def:dyadic-tube}
Let $0<s<1$ and $s<t\leq2$. At sufficiently small dyadic scales,
if $\T$ is a $(\delta,t,\delta^{-\eps})$-set of dyadic tubes and
each tube meets a $(\delta,s,\delta^{-\eps})$-set $\cP(T)$ of
point squares, then
$\#\bigcup_{T\in\T}\cP(T)\geq\delta^{-2s-\eps}$.
\end{theorem}
\begin{proof}
\uses{thm:os-bounded}
Apply point-line duality. Bounded charts and a fixed enlargement
preserve both nonconcentration and incidence, reducing the statement
to the point-to-tube version.
\end{proof}

\section{Elementary incidence bounds}

\begin{proposition}[Incidence bound; OS, Proposition 2.13]
\label{prop:os-incidence}
\uses{def:nice-configuration}
Let $0\leq s\leq t\leq1$, let $\cP$ be a
$(\delta,t,C_P)$-set, and let each $\T(p)$ be a
$(\delta,s,C_T)$-set of $M$ tubes meeting $p$. For any tube family
$\mathcal U$, define
$\mathcal I=\{(p,T):p\in\cP, T\in\mathcal U\cap\T(p)\}$.
Then, with $\theta=(1-t)/(1-s)$,
\[
 \#\mathcal I\lessapprox
 \max\left\{\sqrt{C_PC_T}(M\delta^s)^{\theta/2}
                   (\#\mathcal U)^{1/2}\#\cP, \#\mathcal U\right\}.
\]
For $s=t=1$, set $\theta=0$.
\end{proposition}
\begin{proof}
Count pairs of points on a common tube, split by their separation,
and use the nonconcentration bounds for points and tube directions.
Cauchy--Schwarz converts the resulting pair count into an incidence
bound. The sum over dyadic separations accounts for the logarithmic
loss.
\end{proof}

\begin{corollary}[An elementary tube lower bound; OS, Corollary 2.14]
\label{cor:os-elementary}
\uses{def:nice-configuration}
In the preceding setting, with $\#\T(p)\asymp M$ and
$\T=\bigcup_p\T(p)$,
\[
 \#\T\gtrapprox (C_PC_T)^{-1}M\delta^{-s}
                      (M\delta^s)^{(t-s)/(1-s)}.
\]
At $s=t=1$ the final exponent is interpreted as $1$.
\end{corollary}
\begin{proof}
\uses{prop:os-incidence}
The number of incidences is comparable to $M\#\cP$. Substitute
this into the preceding bound and solve for $\#\T$; the
nonconcentration bounds handle the alternative dominated by
$\#\T$ itself.
\end{proof}

\begin{proposition}[A nonconcentrated coarse tube family; OS, Proposition 4.1]
\label{prop:os-thick-cover}
\lean{quantitative_thick_tube_cover_main}
\uses{def:os-set,def:dyadic-tube}
Let $0<\delta\leq\Delta\leq1$ be dyadic. For each element $p$
of a finite index set $P$, let $\T(p)$ be a
$(\delta,s,C_1)$-family of tubes meeting the unit square, with
$M/2<\#\T(p)\leq M$. There are a subset
$\overline P$ retaining a polylogarithmic fraction of $P$ and a
$(\Delta,s,C_2)$-family $\overline{\T}_\Delta$, where
$C_2\lessapprox_\Delta C_1$, such that each coarse tube contains
at least a constant times
\[
                         H\asymp_\Delta
                         M\#P/\#\overline{\T}_\Delta
\]
incidences $(p,T)$ with $p\in\overline P$ and $T\in\T(p)$.
\end{proposition}
\begin{proof}
Group the fine tubes by their coarse parameter cells, and pigeonhole
the number of incident fine tubes per coarse cell. Apply the fine
nonconcentration estimate to unions of coarse cells. Retaining the
common multiplicity range gives both the coarse Frostman bound and
the lower incidence count.
\end{proof}

\begin{proposition}[Induction on scales; OS, Proposition 5.2]
\label{prop:os-induction}
\uses{def:nice-configuration}
Let $(\cP_0,\T_0)$ be $(\delta,s,C_1,M)$-nice and
$\delta\leq\Delta$ be dyadic. After polylogarithmic refinements,
one obtains a fine configuration, a coarse
$(\Delta,s,C_\Delta,M_\Delta)$-nice configuration on the occupied
$\Delta$-squares, and for each such square $Q$ a rescaled
$(\delta/\Delta,s,C_Q,M_Q)$-nice configuration, with
$C_\Delta,C_Q\asymp_\delta C_1$. These can be chosen so that:
\begin{enumerate}
\item The number of occupied coarse squares and the number of fine
squares in each retained coarse square lose only polylogarithmic factors.
\item Each retained point keeps a polylogarithmic fraction of its
$M$ tubes, and its fine tubes are covered by the coarse tubes over $Q$.
\item The rescaled family preserves the occupied slope intervals
of the fine family at resolution $\delta/\Delta$.
\item The tube counts satisfy
\[
 \frac{\#\T_0}{M}\gtrapprox_\delta
       \frac{\#\T_\Delta}{M_\Delta}
       \max_Q\frac{\#\T_Q}{M_Q}.
\]
\end{enumerate}
\end{proposition}
\begin{proof}
\uses{prop:os-thick-cover}
Regularize the numbers of points, coarse tubes, and fine tubes in
each relevant fibre. The coarse-cover proposition supplies the
coarse nonconcentration. Rescale each retained point square and
track the slope intervals of its fine tubes. Double counting the
regularized incidences gives the product inequality.
\end{proof}

\section{The regular case}

\begin{definition}[Square-root regularity; OS, Definition 6.1]
\label{def:os-regular}
\lean{DirecretisedFurstenbergEstimate.RegularIncidence.IsSquareRootRegular}
\leanok
\uses{def:os-set}
A nonempty $(\delta,t,C)$-set $P$ is
$(\delta,t,C,K)$-regular if
$\cE_{\sqrt\delta}(P)\leq K\delta^{-t/2}$.
Regularity between scales $\delta<\Delta$ means this condition
holds after rescaling every occupied $\Delta$-square to unit size.
\end{definition}

\begin{theorem}[Regular-case incidence gain; OS, Theorem 6.2]
\label{thm:os-regular}
\lean{DirecretisedFurstenbergEstimate.theorem6_1_main}
\leanok
\uses{def:os-regular,def:nice-configuration}
For $0<s<1$ and $s<t<2$, there is $\eps>0$ such that the
following holds at sufficiently small scales. If $P$ is
$(\delta,t,\delta^{-\eps},\delta^{-\eps})$-regular and every
$p\in P$ has a $(\delta,s,\delta^{-\eps})$-family of incident
tubes, then $\cE_\delta(\bigcup_p\T(p))\geq\delta^{-2s-\eps}$.
The threshold can be shared over the larger dimensions used in
the scale-combining argument.
\end{theorem}
\begin{proof}
\uses{prop:os-induction,cor:os-elementary,thm:weak-ring}
\leanok
The sum-product argument gives either an improved fine tube count
or an improved count at scale $\sqrt\delta$. The first alternative
is the conclusion. In the second, apply induction on scales at
$\sqrt\delta$. Square-root regularity keeps a substantial number
of fine points in the retained coarse squares. The elementary
incidence bound in each fibre converts the coarse improvement
into the required fine improvement.
\end{proof}

\section{Combining the scales}

\begin{lemma}[Uniformization; OS, Lemma 7.3]
\label{lem:os-uniformization}
\lean{DiscretisedFurstenbergEstimate.BasicUniformization.basic_uniformization_dyadic}
\uses{def:uniform-set}
For a planar set $P$ and dyadic scales
$1=\Delta_0>\cdots>\Delta_n=\delta$, there is a uniform
refinement $P'$ with
\[
       |P'|_\delta\geq
          \bigl(4n^{-1}\log(1/\delta)\bigr)^{-n}|P|_\delta.
\]
\end{lemma}
\begin{proof}
\uses{lem:uniform-set}
Pigeonhole the branching counts from the bottom scale upwards.
Multiplying the one-level losses and applying the arithmetic-geometric
mean inequality gives the stated dependence on the scale gaps.
\end{proof}

\begin{lemma}[Uniformization after restriction; OS, Lemma 7.4]
\label{lem:os-reuniformization}
\uses{def:uniform-set,def:os-regular}
Suppose $P$ is uniform on $n$ scales and has prescribed
nonconcentration and regularity bounds on specified scale intervals.
If $P'\subset P$ retains at least an $L^{-1}$ fraction of its
finest-scale cells, then $P'$ has a uniform refinement $P''$
retaining at least an $M^{-1}$ fraction of $P$, where
$M=L[4\log(1/\delta)]^n$. The Frostman constants increase by at
most $M$; the square-root covering constants in the regularity
hypotheses do not increase. A union-of-cells structure is preserved.
\end{lemma}
\begin{proof}
\uses{lem:os-uniformization}
Uniformize inside the retained set. The original uniform branching
compares the old and new local normalizations. The loss in total
cardinality weakens a normalized Frostman bound by at most $M$,
while upper covering bounds are inherited by subsets.
\end{proof}

\begin{proposition}[Combining normal, good, and bad scales; OS, Proposition 7.5]
\label{prop:os-combining}
\lean{DiscretisedFurstenbergEstimate.CombiningTheoremRework.prop73_main}
\uses{def:nice-configuration,def:uniform-set,def:os-regular}
Fix $0<s<1$, $t>s$, a scale-separation exponent $\tau>0$, and
a number $n$ of scale blocks. Choose $\eps_G,\eta>0$ sufficiently
small in terms of $s,t$, then $\lambda>0$ sufficiently small in
terms of $s,t,\tau,n$, and $0<\eps_N\leq\eps_G$.
Let a uniform $(\delta,s,\delta^{-\lambda},M)$-nice configuration
have scale ratios $r_j=\Delta_{j-1}/\Delta_j$ and a partition into
normal, good, and bad indices $\mathcal N,\mathcal G,\mathcal B$.
Assume $r_j\geq\delta^{-\tau}$ at normal and good indices.
At a normal index the point set is an
$(s,[\log(1/\delta)]^{C_P}r_j^{\eps_N})$-set between those
scales. At a good index it is regular of some dimension $t_j\geq t$,
with both constants at most
$[\log(1/\delta)]^{C_P}r_j^{\eps_G}$.
Then, at sufficiently small $\delta$,
\[
 \#\T\geq[\log(1/\delta)]^{-C}
       M\delta^{C'\lambda}\delta^{-s+\eps_N}
       \prod_{j\in\mathcal G}r_j^\eta
       \prod_{j\in\mathcal B}r_j^{-1}.
\]
Here $C$ depends on $n,s,\tau,C_P$ and $C'$ on $n,\tau$.
\end{proposition}
\begin{proof}
\uses{prop:os-induction,thm:os-regular,cor:os-elementary,
  lem:os-reuniformization}
Induct on the number of scale blocks. At a normal block use the
elementary estimate; at a good block use the regular-case gain;
at a bad block use the trivial estimate and charge its scale ratio.
After each restriction, uniformize again and pass the resulting
constants to the next block. The product inequality from induction
on scales combines these estimates. The formal contract also
bounds the coefficient of $\lambda$, allowing the final loss to be
chosen before the family is supplied.
\end{proof}

\begin{proposition}[Multiscale decomposition; OS, Proposition 8.1]
\label{prop:os-multiscale}
\lean{DirecretisedFurstenbergEstimate.FinalProof.multiscale_decomposition_main}
\uses{def:uniform-set,def:os-regular}
Let $0<s<1$, $s<t\leq2$, and let $\Delta$ be a sufficiently
small dyadic scale. For a small loss $\eps>0$, a uniform
$(\delta,t,\delta^{-\eps})$-set, $\delta=\Delta^m$ with $m$
large, admits blocks with dimensions $t_j\in[s,2]$, ratios
$r_j=\Delta_{j-1}/\Delta_j$, and structured and bad index sets
$\mathcal S,\mathcal B$ satisfying
\begin{enumerate}
\item $r_j\geq\delta^{-\tau}$ for $j\in\mathcal S$, for some
$0<\tau\leq\eps$, and $\prod_{j\in\mathcal B}r_j\leq\delta^{-\eps}$.
\item On each structured block the set has a $t_j$-Frostman bound
with constant $r_j^\eps$; when $t_j>s$ it is also square-root
regular with that constant.
\item $\prod_{j\in\mathcal S}r_j^{t_j}\geq\delta^{\eps-t}$.
\item Bad blocks are not adjacent.
\end{enumerate}
\end{proposition}
\begin{proof}
\uses{lem:profile-decomposition}
Apply the profile-decomposition lemma to the logarithms of the
uniform branching numbers. Round the resulting interval endpoints
to the permitted grid, discarding short residual intervals.
Superlinear pieces give Frostman bounds, and pieces close to their
chords also give the square-root covering bounds. The total discarded
logarithmic length controls the bad product; the original
$t$-dimensional lower bound gives the structured product.
\end{proof}

\begin{lemma}[Branching profiles and geometry; OS, Lemma 8.3]
\label{lem:profile-geometry}
\uses{def:uniform-set,def:os-regular}
Let $f_P$ be the logarithmic branching profile of a set uniform at
scales $\Delta^j$, $0\leq j\leq m$, with $\delta=\Delta^m$.
If $f_P$ lies above its endpoint chord up to error $\eps m$,
the set has Frostman dimension equal to the slope of that chord,
with constant $O_\Delta(1)\delta^{-\eps}$. If its distance from
the chord is at most $\eps m$, it also has square-root regularity
with the same type of constant.
\end{lemma}
\begin{proof}
Uniformity expresses local cell counts as ratios of products of
branching numbers. Taking logarithms turns the desired inequalities
into comparisons of the profile with its chord.
\end{proof}

\begin{lemma}[Decomposing a branching profile; OS, Lemma 8.5]
\label{lem:profile-decomposition}
Let $0<s<1$, $s<t\leq2$, and $\eps>0$. There is $\tau>0$
such that a $2$-Lipschitz function $f:[0,m]\to\R$ with $f(0)=0$
and $f(x)\geq tx-\eps m$ has disjoint intervals $[c_j,d_j]$
with the following properties. Each has length at least $\tau m$.
On each interval, either $f$ lies within $\eps(d_j-c_j)$ of its
chord and the chord slope is at least $s$, or $f$ lies above that
chord up to the same error and its slope is exactly $s$.
The uncovered length is $O_{s,t}(\eps m)$.
\end{lemma}
\begin{proof}
Use a stopping-time decomposition according to the first slope-$s$
chord that falls below the graph. On the remaining pieces extract
long intervals on which the graph is close to a chord of larger
slope. The Lipschitz bound controls the discarded endpoint pieces,
and the lower bound $tx-\eps m$ prevents them from occupying a
large proportion of the domain.
\end{proof}

\begin{remark}[Loss parameters in Proposition 8.1]
\label{rem:os-losses}
The formalized decomposition keeps a separate bad-scale loss and
the fixed constants from comparing balls and dyadic cells. For
example, the linked interface requires
\[
 \left(1+\frac6{t-s}\right)
 \left(\eps+\frac{2\log9}{\log(1/\Delta)}\right)<t-s.
\]
It returns constants of the form $81\Delta^{-4}r_j^{\eps_{\rm bad}}$.
To obtain the concise paper formulation, choose the source loss and
$\Delta$ sufficiently small before selecting $m$. Identifying
all these losses without making that choice would change the
quantifier order.
\end{remark}
